\documentclass[11pt,a4paper]{article}
\usepackage[utf8]{inputenc}
\usepackage[T1]{fontenc}
\usepackage[portuguese]{babel}
\usepackage{geometry}
\geometry{margin=2.5cm}
\usepackage{listings}
\usepackage{xcolor}
\usepackage{amsmath}
\usepackage{amssymb}
\usepackage{hyperref}
\usepackage{enumitem}
\usepackage{tcolorbox}
\usepackage{tabularx}
\usepackage{booktabs}

\definecolor{codebg}{rgb}{0.95,0.95,0.95}
\definecolor{codegreen}{rgb}{0.1,0.5,0.1}
\definecolor{codepurple}{rgb}{0.5,0.1,0.5}
\definecolor{codegray}{rgb}{0.5,0.5,0.5}

\lstdefinestyle{python}{
    backgroundcolor=\color{codebg},
    commentstyle=\color{codegray}\itshape,
    keywordstyle=\color{codepurple}\bfseries,
    stringstyle=\color{codegreen},
    basicstyle=\ttfamily\small,
    breaklines=true,
    frame=single,
    language=Python,
    showstringspaces=false,
    tabsize=4,
    literate={á}{{\'a}}1 {ã}{{\~a}}1 {é}{{\'e}}1 {í}{{\'i}}1 {ó}{{\'o}}1 {õ}{{\~o}}1 {ú}{{\'u}}1 {ç}{{\c{c}}}1 {Á}{{\'A}}1 {Ã}{{\~A}}1 {É}{{\'E}}1 {Í}{{\'I}}1 {Ó}{{\'O}}1 {Õ}{{\~O}}1 {Ú}{{\'U}}1 {Ç}{{\c{C}}}1 {â}{{\^a}}1 {ê}{{\^e}}1 {ô}{{\^o}}1 {û}{{\^u}}1 {Â}{{\^A}}1 {Ê}{{\^E}}1 {Ô}{{\^O}}1 {Û}{{\^U}}1 {à}{{\`a}}1 {è}{{\`e}}1 {ì}{{\`i}}1 {ò}{{\`o}}1 {ù}{{\`u}}1
}
\lstset{style=python}

\newtcolorbox{objetivos}[1][]{
  colback=green!5!white,
  colframe=green!40!black,
  title=O que você vai aprender nesta página,
  fonttitle=\bfseries,
  #1
}

\newtcolorbox{exercicio}[1]{
  colback=blue!5!white,
  colframe=blue!40!black,
  title=#1,
  fonttitle=\bfseries
}

\newtcolorbox{solucao}{
  colback=yellow!5!white,
  colframe=orange!60!black,
  title=Solução proposta,
  fonttitle=\bfseries
}

\title{\textbf{Data Analysis with Python 2024--2025}\\Parte 3: NumPy (Parte 1)\\\large Soluções dos Exercícios}
\author{Tradução e Adaptação: Universidade de Helsinque (MOOC.fi)}
\date{}

\begin{document}

\maketitle

\section*{Nota Importante}
Este material é uma tradução e adaptação baseada no curso \textit{Data Analysis with Python 2024-2025}, oferecido pela \textbf{Universidade de Helsinque (MOOC.fi)}. As soluções apresentadas a seguir foram elaboradas para fins didáticos, seguindo os enunciados dos exercícios propostos no curso original.

\tableofcontents
\newpage

\section{Soluções - Capítulo 2 (NumPy)}

\subsection*{Exercício 2.11 -- Linhas e colunas}

\begin{solucao}
\begin{lstlisting}
import numpy as np

def get_rows(a):
    return [a[i, :] for i in range(a.shape[0])]

def get_columns(a):
    return [a[:, j] for j in range(a.shape[1])]


def main():
    a = np.array([[5, 0, 3, 3],
                   [7, 9, 3, 5],
                   [2, 4, 7, 6],
                   [8, 8, 1, 6]])
    print(get_rows(a))
    print(get_columns(a))


if __name__ == "__main__":
    main()
\end{lstlisting}

\textbf{Observação:} uma alternativa equivalente, usando a transposição, é implementar \texttt{get\_columns} como \texttt{get\_rows(a.T)}, já que transpor a matriz troca linhas por colunas:

\begin{lstlisting}
def get_columns_alt(a):
    return get_rows(a.T)
\end{lstlisting}
\end{solucao}

\subsection*{Exercício 2.12 -- Vetores-linha e vetores-coluna}

\begin{solucao}
\begin{lstlisting}
import numpy as np

def get_row_vectors(a):
    return [a[i:i+1, :] for i in range(a.shape[0])]

def get_column_vectors(a):
    return [a[:, j:j+1] for j in range(a.shape[1])]


def main():
    a = np.array([[5, 0, 3],
                   [3, 7, 9]])
    print("Vetores-linha:")
    print(get_row_vectors(a))
    print("Vetores-coluna:")
    print(get_column_vectors(a))


if __name__ == "__main__":
    main()
\end{lstlisting}

\textbf{Observação:} o truque usado aqui é indexar com uma fatia de tamanho 1 (\texttt{i:i+1}) em vez de um único índice inteiro (\texttt{i}) -- assim o NumPy preserva aquela dimensão em vez de eliminá-la, o que é exatamente o que diferencia este exercício do anterior.
\end{solucao}

\subsection*{Exercício 2.13 -- Diamante}

\begin{solucao}
\begin{lstlisting}
import numpy as np

def diamond(n):
    # metade superior: identidade seguida da identidade espelhada
    # horizontalmente (concatenadas lado a lado)
    identidade = np.eye(n, dtype=int)
    metade_superior = np.concatenate(
        (identidade, np.fliplr(identidade)),
        axis=1
    )
    # metade inferior: espelho vertical da metade superior,
    # sem repetir a linha central (que ja esta na metade superior)
    metade_inferior = np.flipud(metade_superior)[1:]
    return np.concatenate((metade_superior, metade_inferior))


def main():
    print(diamond(3))
    print(diamond(1))


if __name__ == "__main__":
    main()
\end{lstlisting}
\end{solucao}

\subsection*{Exercício 2.14 -- Comprimento de vetores}

\begin{solucao}
\begin{lstlisting}
import numpy as np

def vector_lengths(a):
    return np.sqrt(np.sum(a**2, axis=1))


def main():
    a = np.array([[3, 4],
                   [1, 0],
                   [0, 0]])
    print(vector_lengths(a))
    # [5. 1. 0.]


if __name__ == "__main__":
    main()
\end{lstlisting}
\end{solucao}

\subsection*{Exercício 2.15 -- Ângulos entre vetores}

\begin{solucao}
\begin{lstlisting}
import numpy as np

def vector_angles(X, Y):
    produto_interno = np.sum(X * Y, axis=1)
    norma_x = np.sqrt(np.sum(X**2, axis=1))
    norma_y = np.sqrt(np.sum(Y**2, axis=1))
    cos_alpha = produto_interno / (norma_x * norma_y)
    cos_alpha = np.clip(cos_alpha, -1.0, 1.0)
    return np.degrees(np.arccos(cos_alpha))


def main():
    X = np.array([[1, 0], [1, 1], [1, 0]])
    Y = np.array([[0, 1], [1, 0], [1, 0]])
    print(vector_angles(X, Y))
    # [90. 45.  0.]


if __name__ == "__main__":
    main()
\end{lstlisting}
\end{solucao}

\subsection*{Exercício 2.16 -- Tabuada revisitada}

\begin{solucao}
\begin{lstlisting}
import numpy as np

def multiplication_table(n):
    linhas = np.arange(n).reshape(n, 1)   # vetor-coluna: 0,1,...,n-1
    colunas = np.arange(n).reshape(1, n)  # vetor-linha: 0,1,...,n-1
    return linhas * colunas               # broadcasting: (n,1) * (1,n) -> (n,n)


def main():
    print(multiplication_table(4))


if __name__ == "__main__":
    main()
\end{lstlisting}
\end{solucao}

\end{document}
